\chapter{Automated Reasoning and Its Limits}
\label{ch:automated-reasoning-limits}

An optimizer proposes replacing an expression by a cheaper one.
Testing can find an input on which the proposal fails, but passing a thousand tests does not establish a claim about every input.
A solver can instead search the space described by a logical formula.
For a local rewrite, the useful query usually asks whether there exists an input satisfying the rewrite's premises on which the two expressions differ.
An unsatisfiable answer rules out such an input, provided the formula faithfully represents the intended semantics.

The engineering work lies on both sides of that answer.
Before the query, choose the domain, encode the instructions, and state the assumptions.
After the query, interpret the result, reconstruct a counterexample when one exists, and preserve a conservative fallback when the solver cannot decide within its limits.
Automated reasoning is most useful when these steps form an explicit contract.

\section{From a theorem to a counterexample query}
\label{sec:ch41-query}

Let $P(x)$ be the precondition of a pure, total rewrite from $s(x)$ to $t(x)$.
The desired theorem is
\[
  \forall x.\ P(x)\rightarrow s(x)=t(x).
\]
Its negation is an existential counterexample claim:
\[
  \exists x.\ P(x)\land s(x)\neq t(x).
\]
A satisfiability solver conventionally searches for values of the free variables in the matrix $P(x)\land s(x)\neq t(x)$; it need not receive an explicit existential quantifier.
The quantifier has not been forgotten: model search supplies its meaning.

For a concrete example, use four-bit unsigned values, with addition reduced modulo $16$.
Consider replacing
\[
 s(x)=\left\lfloor\frac{(x+x)\bmod16}{2}\right\rfloor
 \quad\text{by}\quad t(x)=x.
\]
Without a precondition, $x=8$ refutes the rewrite: the addition wraps to $0$, so $s(8)=0$ while $t(8)=8$.
With $0\leq x<8$, the sum does not wrap, and dividing $2x$ by $2$ gives $x$.
The guarded counterexample query is unsatisfiable.
This is a finite theorem about a particular width, not automatically a theorem about every machine width or source-language integer type.

\term{Boolean satisfiability}, usually abbreviated SAT, asks whether a propositional formula has a satisfying valuation.
\term{Satisfiability modulo theories}, abbreviated SMT, adds interpreted domains and operations such as integer arithmetic, fixed-width bit-vector arithmetic, and array reads and writes.
The distinction matters even for a tiny query.
Treating $x<4$ and $x>9$ as unrelated Boolean atoms permits both to be true; integer arithmetic does not.

An SMT formula also has a \term{sort} for each term, such as Boolean, integer, or a bit vector of a particular width.
Sort checking catches malformed combinations, but it cannot detect that the encoder chose integers when the instruction actually wraps.
A well-sorted formula can still prove the wrong proposition.

\section{Boolean search and its representation}
\label{sec:ch41-sat}

A \term{literal} is an atom or its negation.
A \term{clause} is a disjunction of literals, and \term{conjunctive normal form}, or CNF, is a conjunction of clauses.
The empty clause is false; an empty collection of clauses is true.
The formula
\[
 (a\lor b)\land(\neg a\lor b)\land(\neg b\lor c)\land\neg c
\]
is unsatisfiable.
The unit clause $\neg c$ forces $c$ false.
Then $\neg b\lor c$ forces $b$ false, after which the first two clauses demand both $a$ and $\neg a$.
This sequence illustrates \term{unit propagation}: if every literal but one in a clause is false, the remaining literal must be true.
An empty residual clause exposes a conflict.

When propagation does not settle the formula, a complete search can choose an unassigned atom, try one truth value, and backtrack to the other when necessary.
For $k$ atoms there are $2^k$ complete valuations.
Evaluating a formula with $m$ literal occurrences under every valuation takes $O(2^k m)$ time with a straightforward representation.
That bound describes a baseline procedure, not the behavior of a modern SAT solver on every input.
Clause learning, better branching choices, and efficient propagation can avoid large parts of this search.
They do not turn arbitrary SAT instances into a known polynomial-time decision problem; general SAT is NP-complete~\cite{huth-2004-logic}.

The representation delivered to the solver affects both correctness and size.
Converting formulas to CNF solely by distributing disjunction over conjunction can cause exponential expansion.
For example, distributing a disjunction of $k$ separate two-literal conjunctions produces $2^k$ clauses in the direct expansion.
A compiler expression already has a graph of intermediate operations.
Naming those operations in the formula avoids repeatedly expanding their definitions.

\subsection{Tseitin encoding preserves satisfiability}
\label{sec:ch41-tseitin}

A \term{Tseitin encoding} introduces an auxiliary Boolean variable for each internal formula node and constrains it to equal that node's operation.
For a conjunction node $z\leftrightarrow(a\land b)$, use
\[
 (\neg z\lor a)\land(\neg z\lor b)\land
 (z\lor\neg a\lor\neg b).
\]
The first two clauses say that a true $z$ requires both inputs true.
The third says that two true inputs require $z$ true.
Both directions are necessary for this full definitional encoding.
Table~\ref{tab:ch41-gate-clauses} gives the corresponding fixed-size encodings for the operators used below.

\begin{table}
\caption{Each gate is constrained by clauses equivalent to its local definition.}
\label{tab:ch41-gate-clauses}
\centering
\begin{tabular}{p{0.25\linewidth}p{0.64\linewidth}}
\toprule
Definition & Clauses, all conjoined \\
\midrule
$z\leftrightarrow\neg a$ & $(\neg z\lor\neg a)$; $(z\lor a)$ \\
$z\leftrightarrow(a\land b)$ & $(\neg z\lor a)$; $(\neg z\lor b)$; $(z\lor\neg a\lor\neg b)$ \\
$z\leftrightarrow(a\lor b)$ & $(z\lor\neg a)$; $(z\lor\neg b)$; $(\neg z\lor a\lor b)$ \\
\bottomrule
\end{tabular}
\end{table}

Finally, assert the root's variable with a unit clause.
Without that assertion, the definitions merely describe evaluation: every assignment to the original atoms extends to a consistent assignment of gate values, whether the original formula is true or false.
Algorithm~\ref{alg:ch41-encode} constructs the full encoding for a finite formula DAG.
Inputs that are already literals need no auxiliary node, so the gate clauses may refer directly to a negated atom.

\begin{algorithm}
\caption{Encode a Boolean formula DAG as a satisfiability-equivalent CNF.}
\label{alg:ch41-encode}
\begin{pseudocode}
\Require Root $r$ of an acyclic DAG with literal leaves and unary NOT or binary AND/OR nodes
\Ensure Clauses satisfiable exactly when the formula rooted at $r$ is satisfiable
\State $K\gets$ empty clause list
\ForAll{nodes $v$ in a children-before-parent topological order}
  \If{$v$ is a literal}
    \State $L[v]\gets$ that literal
  \Else
    \State $z\gets$ a fresh atom; $L[v]\gets z$
    \If{$v$ is NOT with child $a$}
      \State Append $(\neg z\lor\neg L[a])$ and $(z\lor L[a])$ to $K$
    \ElsIf{$v$ is AND with children $a,b$}
      \State Append $(\neg z\lor L[a])$ and $(\neg z\lor L[b])$ to $K$
      \State Append $(z\lor\neg L[a]\lor\neg L[b])$ to $K$
    \Else \Comment{OR with children $a,b$}
      \State Append $(z\lor\neg L[a])$ and $(z\lor\neg L[b])$ to $K$
      \State Append $(\neg z\lor L[a]\lor L[b])$ to $K$
    \EndIf
  \EndIf
\EndFor
\State Append unit clause $(L[r])$ to $K$
\State \Return $K$
\end{pseudocode}
\end{algorithm}

Let $N$ be the number of distinct DAG nodes and $E$ its edge count.
Topological traversal takes $O(N+E)$ time and space; bounded fan-in gives $E=O(N)$.
Each internal node adds one atom and at most three clauses, each of length at most three.
With constant-time allocation of identifiers and literals, encoding therefore takes $O(N)$ time and space.
This bound concerns construction only, not solving.
Shared nodes must be visited once; expanding the DAG into a tree first can destroy the bound relative to the original representation.

To see the semantic guarantee, let $D(\mathbf p,\mathbf z)$ be all gate-definition clauses, $\mathbf p$ the original atoms, and $z_r$ the root atom.
For every fixed valuation of $\mathbf p$, structural evaluation determines exactly one valuation of $\mathbf z$ satisfying $D$.
In that extension, $z_r$ equals the original formula $F(\mathbf p)$.
Consequently,
\[
 F(\mathbf p)\quad\Longleftrightarrow\quad
 \exists\mathbf z.\,D(\mathbf p,\mathbf z)\land z_r.
\]
The formulas are \term{equisatisfiable}: either both have a model or neither does.
They are not interchangeable as unquantified formulas over the enlarged vocabulary, because an arbitrary assignment to the auxiliary atoms need not satisfy the definitions.
The stronger projection property displayed above is what lets the checker recover a counterexample in the original inputs.

\subsection{A complete small encoding}
\label{sec:ch41-encoding-example}

Take $F=(p\land q)\lor(\neg p\land r)$.
Introduce $a\leftrightarrow(p\land q)$, $b\leftrightarrow(\neg p\land r)$, and $c\leftrightarrow(a\lor b)$, as in Figure~\ref{fig:ch41-gates}.
There are three gate definitions, nine definition clauses, and the unit clause $(c)$: ten clauses in total over six atoms.
Negated $p$ is a literal, so it needs no separate NOT gate in this example.

\begin{figure}
\centering
\begin{tikzpicture}[node distance=13mm and 17mm]
\node[draw, rounded corners, align=center] (a) {$a\leftrightarrow p\land q$};
\node[draw, rounded corners, align=center, right=of a] (b) {$b\leftrightarrow\neg p\land r$};
\node[draw, rounded corners, align=center, below=of a, xshift=22mm] (c) {$c\leftrightarrow a\lor b$};
\node[draw, double, below=of c] (root) {assert $c$};
\draw[cfg edge] (a) -- (c);
\draw[cfg edge] (b) -- (c);
\draw[cfg edge] (c) -- (root);
\end{tikzpicture}

\caption{Auxiliary atoms name intermediate Boolean values, and the root is separately asserted.}
\label{fig:ch41-gates}
\end{figure}

Table~\ref{tab:ch41-gate-values} evaluates every input valuation.
Exactly four of the eight input rows satisfy $F$.
Each of those rows has exactly one extension satisfying the full ten-clause encoding, so exactly four of the $64$ assignments to all six atoms satisfy it.
For $p=0,q=1,r=1$, the extension is $a=0,b=1,c=1$.
The satisfying model projected onto $p,q,r$ is the original witness.

\begin{table}
\caption{Evaluating every gate fixes a unique extension of each input valuation.}
\label{tab:ch41-gate-values}
\centering
\begin{tabular}{cccccc}
\toprule
$p$ & $q$ & $r$ & $a=p\land q$ & $b=\neg p\land r$ & $c=a\lor b$ \\
\midrule
0 & 0 & 0 & 0 & 0 & 0 \\
0 & 0 & 1 & 0 & 1 & 1 \\
0 & 1 & 0 & 0 & 0 & 0 \\
0 & 1 & 1 & 0 & 1 & 1 \\
1 & 0 & 0 & 0 & 0 & 0 \\
1 & 0 & 1 & 0 & 0 & 0 \\
1 & 1 & 0 & 1 & 0 & 1 \\
1 & 1 & 1 & 1 & 0 & 1 \\
\bottomrule
\end{tabular}
\end{table}

If the root clause is omitted, all eight input valuations have a consistent extension, including those where $F$ is false.
If only $z\rightarrow(a\land b)$ is encoded for a conjunction, false gate values can hide true operands.
Such one-directional constraints can be appropriate in specialized polarity-aware encodings, but they need their own correctness argument.
They are not a drop-in replacement for the full definitions in Algorithm~\ref{alg:ch41-encode}.

\section{The theory is part of the specification}
\label{sec:ch41-theories}

Over mathematical integers, $x+3>x$ holds for every $x$.
For four-bit unsigned values, the addition wraps, and the same-looking comparison fails near the upper bound.
Table~\ref{tab:ch41-wrap} gives the boundary cases.
The rewrite to true is valid under $x\leq12$ and invalid without that premise.

\begin{table}
\caption{Four-bit addition followed by unsigned comparison disagrees with integer arithmetic after wraparound.}
\label{tab:ch41-wrap}
\centering
\begin{tabular}{rrc}
\toprule
$x$ & $(x+3)\bmod16$ & $((x+3)\bmod16)>x$ \\
\midrule
0 & 3 & true \\
12 & 15 & true \\
13 & 0 & false \\
14 & 1 & false \\
15 & 2 & false \\
\bottomrule
\end{tabular}
\end{table}

A width-$w$ \term{bit vector} has $2^w$ possible bit patterns.
Modular addition is defined on those patterns, but the ordering operation must still be chosen.
The pattern $1111$ represents $15$ under unsigned interpretation and $-1$ under four-bit two's-complement signed interpretation.
Thus the comparison of that pattern with $0000$ depends on whether the comparison is unsigned or signed.
The addition bits alone do not determine that choice.

Signed source arithmetic may impose a different definedness rule from raw bit-vector addition.
For example, a four-bit signed model with an overflow fault cannot execute $7+1$ normally, whereas modular bit-vector addition produces the pattern for $-8$.
Encoding only the modular result would omit the fault condition.
For C, unsigned arithmetic and signed overflow have different rules, and narrow integer operands can be promoted before arithmetic~\cite{wg14-2011-c11}.
An eight-bit storage type therefore does not by itself imply an eight-bit addition in a source expression.
The checker must encode the actual operation and conversions.

Integer theories can be useful when analysis proves that a machine operation stays within range.
One can then prove an integer identity together with obligations that every relevant intermediate result avoids overflow.
Dropping the range obligations changes the theorem.
The same care applies to division: mathematical integer division, truncation toward zero, and unsigned division must not be conflated for negative operands.
Fixed-width division can also have exceptional source or IR cases even when the solver's corresponding operator is total.

To express such cases, associate a definedness predicate $D_e(x)$ with an expression and a value expression $V_e(x)$ for its normal result.
For a deterministic replacement that must preserve normal behavior whenever the source is defined, search for
\[
 P(x)\land D_s(x)\land
 \bigl(\neg D_t(x)\lor V_s(x)\neq V_t(x)\bigr).
\]
This formula assumes that normal values are the only observations being compared and that the value encodings are defined as logical terms even outside their meaningful program domains.
If traps are observable and must be preserved, compare the relevant outcomes explicitly instead; source trapping inputs cannot simply be excluded.
Memory, nondeterminism, and concurrency need richer behavior relations than this scalar formula.

SMT is a family of problems, not a promise that every combination of operations is easy or decidable.
Quantifier-free fixed-width bit-vector formulas have finite domains and are decidable, often by translating operations into Boolean circuits and then SAT.
Linear integer arithmetic also has decision procedures, despite its infinite domain.
By contrast, unrestricted integer multiplication brings in much harder questions, and arbitrary quantifiers or theory combinations require their own analysis.
Choosing a solver's logic mode should follow the actual encoded operations, not the broad label \enquote{arithmetic}.

\section{Verification conditions connect execution to formulas}
\label{sec:ch41-vcs}

A \term{verification condition} is a logical obligation whose validity suffices for a stated program property.
The weakest-precondition construction of Chapter~\ref{ch:logic-compiler-reasoning} produces one such obligation: $P\rightarrow\operatorname{wp}(C,Q)$.
Checking its validity can be reduced to checking unsatisfiability of $P\land\neg\operatorname{wp}(C,Q)$.
For larger expressions, forward symbolic execution with fresh logical names can avoid duplicating substituted terms.
Each name denotes a particular value at a particular point, even when the source variable is assigned repeatedly.

Revisit the integer conditional that negates negative $x$, adds $2$ otherwise, and finally sets $y$ to three times the resulting $x$.
Use $x_0$ for the initial value and $x_1,y_1$ for the resulting values.
The transition relation is
\[
\begin{split}
 T(x_0,x_1,y_1)={}&
 \bigl((x_0<0\land x_1=-x_0)\lor
       (x_0\geq0\land x_1=x_0+2)\bigr)\\
 &{}\land y_1=3x_1.
\end{split}
\]
The names are logical bookkeeping; no assignment can retroactively change $x_0$.
For precondition $x_0\neq-1$ and goal $y_1\geq6$, the counterexample query is
\[
  x_0\neq-1\land T(x_0,x_1,y_1)\land y_1<6.
\]
It is unsatisfiable over integers.
On the negative arm, $x_0\neq-1$ forces $x_0\leq-2$, so $x_1\geq2$.
On the other arm, $x_0\geq0$ also gives $x_1\geq2$.
Removing the precondition yields the model $x_0=-1,x_1=1,y_1=3$.
Thus both the backward and forward encodings expose the same missing premise.

An encoder must represent every behavior relevant to the claim.
An underapproximation that leaves out the negative branch can make this incorrect unconditional claim appear proved.
An overapproximation that includes extra behaviors is safer for proving absence of counterexamples, but it can produce spurious models that do not correspond to executions.
Replaying a model distinguishes a real transformation bug from an abstraction artifact or an encoding error.

For a loop invariant $I$, generate separate obligations for initialization, preservation, and exit.
Using a transition relation $T_C(\rho,\rho')$ for one body execution, the preservation counterexample query has the form
\[
  I(\rho)\land b(\rho)\land T_C(\rho,\rho')\land\neg I(\rho').
\]
An unsatisfiable result proves closure of the proposed invariant under the body; it does not invent an invariant that proves the desired exit property.
Nor does it prove termination.
A termination proof adds a well-founded variant and obligations that each taken iteration decreases it and that the body terminates.

\term{Bounded model checking} instead unrolls a transition relation a finite number of times.
With bound $k$, an unsatisfiable query can rule out the encoded bad executions within that bound.
It does not rule out a failure at iteration $k+1$ unless a completeness bound or an inductive argument closes the gap.
For example, a counter initialized to zero that increments until $10$ and then violates an assertion has no failure within its first nine increments.
The bound is part of the theorem and belongs in the report.

\section{Solver results and the trusted boundary}
\label{sec:ch41-results}

Figure~\ref{fig:ch41-checker-pipeline} shows a practical local-checking path.
The semantic description and precondition feed the encoder; the solver receives only the resulting formula.
Table~\ref{tab:ch41-outcomes} states what each solver outcome supports.

\begin{figure}
\centering
\begin{tikzpicture}[node distance=10mm]
\node[draw, rounded corners, align=center, text width=65mm] (rule) {typed source, replacement, and premises};
\node[draw, rounded corners, align=center, text width=65mm, below=of rule] (query) {encode a defined counterexample};
\node[draw, rounded corners, align=center, text width=65mm, below=of query] (solver) {solver: SAT, UNSAT, or UNKNOWN};
\node[draw, rounded corners, align=center, text width=65mm, below=of solver] (result) {replay witness, validate theorem, or retain original};
\draw[cfg edge] (rule) -- (query);
\draw[cfg edge] (query) -- (solver);
\draw[cfg edge] (solver) -- (result);
\end{tikzpicture}

\caption{A solver result becomes evidence about a rewrite only through a faithful encoding and result interpretation.}
\label{fig:ch41-checker-pipeline}
\end{figure}

\begin{table}
\caption{Counterexample-query outcomes have different consequences for an optimization.}
\label{tab:ch41-outcomes}
\centering
\begin{tabular}{p{0.15\linewidth}p{0.36\linewidth}p{0.36\linewidth}}
\toprule
Result & Meaning for the formula & Action for the proposed rewrite \\
\midrule
SAT & At least one satisfying assignment exists. & Project the model onto inputs and replay the proposed counterexample. \\
UNSAT & No assignment satisfies the query in the selected theory. & Accept the conditional theorem if encoding and premises are justified. \\
UNKNOWN & The solver has established neither satisfiability nor unsatisfiability. & Retain the original program or use an independently justified rule. \\
\bottomrule
\end{tabular}
\end{table}

An interface may report resource exhaustion or cancellation separately from UNKNOWN.
For optimization decisions, neither supports acceptance of an otherwise unproved rewrite.
UNKNOWN is not a third semantic truth value of the formula.
A satisfiable formula remains satisfiable even if a particular run cannot find its model.
A model supplied after an inconclusive result should not be treated as a counterexample until independently checked against the query.

A SAT result is especially useful during rule development.
Decode widths and signedness, evaluate the original and replacement IR on the input, and display the first differing observation.
For the four-bit doubling example, report $x=8$, wrapped sum $0$, source result $0$, and replacement result $8$.
Reporting only six auxiliary Boolean values from a circuit encoding would obscure the bug.
Retain the model together with the exact precondition and type parameters as a regression case.

An UNSAT result can be accompanied by a proof certificate when the solver and chosen mode support one.
A separate checker can validate that certificate and reduce reliance on a large search implementation.
It still checks a theorem about the encoded formula.
The instruction semantics, query construction, certificate checker, and application of the resulting rewrite remain part of the trusted path.
Certificate checking is therefore valuable without being a substitute for a semantic specification.

For repeated queries, shared expression nodes and an incremental solver context can reuse common premises.
Scope each query's extra assumptions carefully: an assumption leaked from a previous rewrite can make a later query vacuously unsatisfiable.
Cache keys must include operator semantics, widths, signedness, overflow conditions, and the premises under which the result was established.
A theorem under $x<8$ is not an unconditional theorem merely because the expression graph matches a previous query.

\section{What undecidability rules out}
\label{sec:ch41-limits}

There are two distinct reasons a checker may fail to answer promptly.
A decidable problem can require too much time or memory on a particular input.
An undecidable problem admits no algorithm that terminates with a correct yes-or-no answer on every input in its specified class.
The second is a limit on every possible algorithm, not a complaint about the performance of the current implementation~\cite{boolos-2007-computability}.

Arbitrary first-order validity is undecidable, although standard first-order calculi are sound and complete.
There is no contradiction.
An effective enumeration of finite proofs eventually finds a proof of each valid sentence, but need not terminate on an invalid sentence.
Completeness guarantees existence of a finite proof when validity holds; it does not give a terminating unsuccessful-search test.
The result concerns validity across first-order structures and must not be confused with deciding a particular restricted theory or evaluating a formula in one finite structure.

A similar boundary appears for whole-program equivalence in a language capable of unbounded computation.
Given a machine $M$ and an input $u$, construct a program $A$ that simulates $M(u)$ and returns $0$ if the simulation halts.
Construct $B$ that loops forever.
Assume the simulation has no externally visible effects and that termination is observable.
Then $A$ and $B$ have the same behavior exactly when $M(u)$ does not halt.
A total equivalence decider for all such programs would therefore decide the halting problem by complementing its answer.
No such decider exists~\cite{boolos-2007-computability}.
The reduction states its computational and observational assumptions; it is not a claim that every local optimization problem is undecidable.

Fixed-width straight-line arithmetic with finitely many inputs has a finite input space.
Exhaustive enumeration is a decision procedure, even if the number of assignments is prohibitive.
A fixed finite-state machine also has decidable reachability: explore reachable states until every successor has been visited or a bad state is found.
For a deterministic machine with at most $2^m$ complete states, revisiting a state repeats its future behavior, so an unbounded run can be detected after a finite exploration.
The state must include everything that determines the next step, including control location and any modeled environment state.
Unbounded memory, unbounded integers, or an unmodeled external environment invalidate that finite-state argument.

These distinctions suggest a workable compiler design.
Use sound analysis to obtain premises; restrict local queries to theories suited to the operations; use invariants or summaries at loops and calls; and apply time limits with conservative fallbacks.
Check a single candidate transformation rather than claiming to decide equivalence for every possible program.
That restricted task can have a complete finite decision procedure even though unrestricted program equivalence cannot.
The limit tells us how to state the service contract, not that automation is futile.

\section{Exercises}
\label{sec:ch41-exercises}

Solutions and C11 checks are in \path{example-code/logic-checks/}.

\begin{enumerate}
\item\label{ex:ch41-root}
Encode $p\land\neg p$ using a fresh conjunction variable $z$ and the clauses of Table~\ref{tab:ch41-gate-clauses}.
Show that the encoding with unit clause $(z)$ is unsatisfiable.
Give a model after removing that unit clause.

\item\label{ex:ch41-projection}
For the six-atom encoding in Section~\ref{sec:ch41-encoding-example}, count the satisfying assignments with and without the root clause.
Explain why the count does not multiply by $2^3$ when auxiliary atoms are introduced.

\item\label{ex:ch41-overflow}
For four-bit unsigned addition, give the weakest condition on $x$ under which $((x+3)\bmod16)>x$ is true.
List all counterexamples when the condition is absent.

\item\label{ex:ch41-width}
Generalize the doubling-and-halving rewrite of Section~\ref{sec:ch41-query} to width $w\geq1$.
Derive the exact subset of inputs on which it equals $x$, and a formula for the erroneous result on all remaining inputs.

\item\label{ex:ch41-vc}
For the transition relation in Section~\ref{sec:ch41-vcs}, change the postcondition to $y_1\geq9$.
Find the weakest condition on integer $x_0$ that makes the counterexample query unsatisfiable, and give every excluded integer input.

\item\label{ex:ch41-unknown}
A rewrite checker accepts a rule whenever the solver does not return SAT.
Explain the error, repair the acceptance condition, and state what an UNSAT answer at an unrolling bound of $20$ proves when no completeness argument is available.
\end{enumerate}

\section{Further reading}
\label{sec:ch41-reading}

Huth and Ryan~\cite{huth-2004-logic} provide foundations for satisfiability, proof systems, and verification conditions, including the distinction between semantic claims and procedures that search for evidence.
Boolos, Burgess, and Jeffrey~\cite{boolos-2007-computability} develop reductions, first-order undecidability, and the relation between completeness and effective proof search; these results explain exactly which universal checker specifications are impossible.
The C11 committee draft~\cite{wg14-2011-c11}, particularly its integer promotions and expression rules, is a useful exercise in translating a language specification into explicit arithmetic and definedness premises.
